\clearpage
\SourceLeaf{70}
\pdfbookmark[1]{Observationes cometae secundi a. 1813. Review of 1814}{anzeige-cometae1814}
\begin{center}\rule{.75\textwidth}{.3pt}\\
Göttingische gelehrte Anzeigen. 3 January 1814.\\
\rule{.75\textwidth}{.3pt}\end{center}
We still owe a review of a short paper that Professor \textsc{Gauss} presented to the Royal Society on 10 September of last year, entitled
\begin{center}
\textit{Observationes cometae secundi a.} 1813 \textit{in observatorio Gottingensi factae, adiectis
nonnullis adnotationibus circa calculum orbitarum parabolicarum.}
\end{center}
As is well known, it was our Professor \textsc{Harding} who discovered this comet on 3 April of last year in Poniatowski’s Bull. The observations made at our observatory run from 9 to 25 April. They are given \textit{in full} in the paper. Since most have already been published elsewhere, we pass over them here in silence and give only the parabolic orbit that Professor \textsc{Gauss} asked Dr. \textsc{Gerling} in Cassel to calculate from our observations and some from Bremen and Paris.
\begin{center}\small
Time of passage through perihelion, 19 May 1813, $10^{\mathrm h}\;24^{\mathrm m}\;5^{\mathrm s}$\\
Göttingen mean time
\end{center}
\begin{center}\small
\setlength{\tabcolsep}{3pt}
\begin{tabular}{ll}
Longitude of perihelion $\quad.\;.\;.\;.$ & $197^0\;43'\;7''7$\\
Longitude of the ascending node $\quad,$ & $42\;40\;15.2$\\
Orbital inclination $\quad.\;.\;.\;.\;.\;.$ & $81\quad2\;11.8$\\
Logarithm of the least distance & $0.0849212$\\
\multicolumn{2}{l}{Motion \textit{retrograde.}}
\end{tabular}\end{center}

\SourceLeaf{71}%
Only a few decades ago, few people in all of Europe could calculate a comet’s orbit. Improved methods have now made the work so much easier and simpler that an otherwise capable person can learn it without difficulty and determine a comet’s orbit in fewer hours than the number of days formerly required. For the \textit{first} calculation of an orbit from three observations, Dr. \textsc{Olbers}’s convenient method leaves almost nothing to be desired, provided one can choose observations for which the direction of the geocentric motion does not come too close to the direction of the great circle drawn on the celestial sphere through the middle geocentric position of the comet and the corresponding position of the Sun. The second part of the present paper aims to simplify a few points of this elegant method still further and make them still more convenient for numerical calculation. The attempt concerns the indirect procedure for determining approximate values of the comet’s distances from the Sun and of the chord between the two outermost positions. Only someone inexperienced in astronomical calculation could fault the method for being indirect. But the form that \textsc{Olbers} gave it has this inconvenience. In most cases the smaller five-decimal tables would yield errors that are too large even if, as is reasonable for a first approximation, one does not seek the utmost precision in the calculation. This is because several quantities that are very small by nature occur as differences of other quantities that become nearly equal. Professor \textsc{Gauss} sought to meet this difficulty by introducing several auxiliary quantities that can always be determined with sufficient accuracy from the small tables. He thereby gave the formulas for the distances from the Sun and for the chord a different form that also seems to make the subsequent trials of the indirect solution somewhat easier. These transformations do not lend themselves to an extract unless we are willing to copy this part of the paper out in full. We therefore refer readers to the second volume of the Society’s \textit{Commentationes}, which is to appear shortly and in which the entire paper is already printed.

The author thought it would be a service to some readers if, alongside his own changes to the method, he provided
